% from ega2-1-fr.tex
\oldpage[II]{131}

le noyau de cet homomorphisme, qui est un Idéal cohérent de
$g_*(\mathcal{O}_{X'})=\mathcal{B}$
(\hyperref[I.6.1.1-fr]{I, 6.1.1}), et posons
$X''=\operatorname{Spec}(\mathcal{B}/\mathcal{T})$~; on a donc un diagramme
commutatif
\[
\xymatrix{
X''\ar[rr]^h\ar[dr]_{g'} && X'\ar[dl]^g\\
& X
}
\]
où $h$ est une immersion fermée (\hyperref[II.1.4.10-fr]{1.4.10}). En outre,
on sait que le support de $\mathcal{T}$ est un ensemble fermé
(\hyperref[0.5.2.2-fr]{0, 5.2.2}) rare dans $X$
(\hyperref[I.7.4.6-fr]{I, 7.4.6}), d'où on conclut que pour le point générique
$x$ d'une composante irréductible de $X$, il y a un voisinage ouvert affine
$U$ de $x$ tel que
$\mathcal{B}|U=(\mathcal{B}/\mathcal{T})|U$. Comme $g$ est par hypothèse
surjectif, on en conclut que $x\in g'(X'')$~; $g'$ est donc dominant, et étant
un morphisme fini, il est \emph{surjectif}
(\hyperref[II.6.1.10-fr]{6.1.10})~; on a par définition
\[
g'_*(\mathcal{O}_{X''})\otimes_{\mathcal{O}_X}\mathcal{R}(X)
=(\mathcal{B}/\mathcal{T})\otimes_{\mathcal{O}_X}\mathcal{R}(X)
=g_*(\mathcal{O}_{X'})\otimes_{\mathcal{O}_X}\mathcal{R}(X),
\]
donc $(X,\mathcal{O}_X)$ et $g'_*(\mathcal{O}_{X''})$ vérifient
(II \emph{bis}), et en outre
\[
g'_*(\mathcal{O}_{X''})\longrightarrow
g'_*(\mathcal{O}_{X''})\otimes_{\mathcal{O}_X}\mathcal{R}(X)
\]
est injectif. Enfin, $h^*(\mathcal{L}')=\mathcal{L}''$ est un
$\mathcal{O}_{X''}$-Module ample
(\hyperref[II.4.6.13-fr]{4.6.13, (i \emph{bis})}), et on a
$N_{X''/X}(\mathcal{L}'')=N_{X'/X}(\mathcal{L}')$. En effet, pour définir
ces deux $\mathcal{O}_X$-Modules inversibles on peut utiliser un même
recouvrement ouvert affine $(U_\lambda)$ de $X$ tel que les restrictions de
$g_*(\mathcal{L}')$ et $g'_*(\mathcal{L}'')$ à $U_\lambda$ soient
respectivement isomorphes à $\mathcal{B}|U_\lambda$ et
$(\mathcal{B}/\mathcal{T})|U_\lambda$. On voit aussitôt qu'à tout
isomorphisme
\[
\eta_\lambda:g_*(\mathcal{L}')|U_\lambda\longrightarrow
\mathcal{B}|U_\lambda
\]
correspond canoniquement un isomorphisme
\[
\eta'_\lambda:g'_*(\mathcal{L}'')|U_\lambda\longrightarrow
(\mathcal{B}/\mathcal{T})|U_\lambda,
\]
de sorte que, si $(\omega_{\lambda\mu})$ et
$(\omega'_{\lambda\mu})$ sont les $1$-cocycles correspondant aux systèmes
d'isomorphismes $(\eta_\lambda)$ et $(\eta'_\lambda)$
(\hyperref[II.6.5.2-fr]{6.5.2}), $\omega'_{\lambda\mu}$ soit l'image
canonique dans
$\Gamma(U_\lambda\cap U_\mu,\mathcal{B}/\mathcal{T})$ de
$\omega_{\lambda\mu}\in\Gamma(U_\lambda\cap U_\mu,\mathcal{B})$. En vertu
de la définition de $\mathcal{T}$, on en conclut que
\[
N_{\mathcal{B}/\mathcal{A}}\omega_{\lambda\mu}
=N_{(\mathcal{B}/\mathcal{T})/\mathcal{A}}\omega'_{\lambda\mu}
\]
(avec $\mathcal{A}=\mathcal{O}_X$), d'où l'égalité énoncée.

Supposons donc l'homomorphisme
$g_*(\mathcal{O}_{X'})\to
g_*(\mathcal{O}_{X'})\otimes_{\mathcal{O}_X}\mathcal{R}(X)$ injectif
lorsqu'on est dans l'hypothèse (II \emph{bis}). Il suffit de prouver que
lorsque $f$ parcourt les sections de $\mathcal{L}^{\otimes n}$ ($n>0$)
au-dessus de $X$, les $X_f$ forment une base de la topologie de $X$
(\hyperref[II.4.5.2-fr]{4.5.2}). Or, soit $x\in X$, et soit $U$ un voisinage
quelconque de $x$~; comme $g^{-1}(x)$ est fini
(\hyperref[II.6.1.7-fr]{6.1.7}) et que $\mathcal{L}'$ est ample, il existe un
entier $n>0$ et une section $f'$ de ${\mathcal{L}'}^{\otimes n}$ au-dessus de
$X'$, tels que $X'_{f'}$ soit un voisinage de $g^{-1}(x)$ contenu dans
$g^{-1}(U)$ (\hyperref[II.4.5.4-fr]{4.5.4}). Comme
\[
\mathcal{L}^{\otimes n}=N_{X'/X}({\mathcal{L}'}^{\otimes n}),
\]
il suffit de prendre $f=N_{X'/X}(f')$~; en effet, alors
$X-X_f=g(X'-X'_{f'})$ (\hyperref[II.6.5.7-fr]{6.5.7}), donc
$x\in X_f\subset U$.

\begin{corollary}[6.6.3]
\label{II.6.6.3-fr}
Sous les hypothèses de (\hyperref[II.6.6.1-fr]{6.6.1}), pour qu'un
$\mathcal{O}_X$-Module inversible $\mathcal{L}$ soit ample pour $f$, il faut
et il suffit que $\mathcal{L}'=g^*(\mathcal{L})$ soit ample pour $f\circ g$.
\end{corollary}

La condition est nécessaire, puisque $g$ est affine
(\hyperref[II.5.1.12-fr]{5.1.12}). Pour démontrer que la condition est
suffisante, on peut supposer $Y$ affine (\hyperref[II.4.6.4-fr]{4.6.4}), donc
$X$ et $X'$ quasi-compacts et $\mathcal{L}'$ ample
(\hyperref[II.4.6.6-fr]{4.6.6}), et il faut montrer que $\mathcal{L}$ est
ample. Or, l'ensemble des points

